\worldchapter{\trans{weapon_modifications_title}}{appendix-weaponmodifications}
\label{ch:appendix-weaponmodifications}

\begin{multicols}{2}

    \section{Lauf}



    \subsection*{Mittlerer Schalldämpfer}

    \begin{pexample}
    Der Mittlere Schalldämpfer reduziert den Schussknall erheblich, verringert aber auch die Durchschlagskraft an der Waffe.
    \end{pexample}


    \begin{tabular}{@{}ll@{}}
        \textbf{\trans{price_label}:} & 200 \\
        \textbf{Verborgenheit} & -2\\
        \textbf{Durchschlag} & -1\\
    \end{tabular}

    \wbif{phasesix}{
        \vspace{3mm}
\faIcon{flag}\hspace{0.5em}\faIcon{plane}\hspace{0.5em}\faIcon{microchip}\hspace{0.5em}\faIcon{space-shuttle}\hspace{0.5em}    }{}



    \subsection*{Low-Profile Schalldämpfer}

    \begin{pexample}
    Der Low-Profile-Schalldämpfer reduziert den Schussknall erheblich, verringert aber auch den Schaden an der Waffe.
    \end{pexample}


    \begin{tabular}{@{}ll@{}}
        \textbf{\trans{price_label}:} & 400 \\
        \textbf{Schadenspotential} & -1\\
        \textbf{Verborgenheit} & -2\\
    \end{tabular}

    \wbif{phasesix}{
        \vspace{3mm}
\faIcon{plane}\hspace{0.5em}\faIcon{microchip}\hspace{0.5em}    }{}


    \section{Visiere}



    \subsection*{Holo Visier}

    \begin{pexample}
    Das Holovisier erhöht die Genauigkeit und das Schadenspotential der Waffe auf mittlere Distanzen.
    \end{pexample}


    \begin{tabular}{@{}ll@{}}
        \textbf{\trans{price_label}:} & 1800 \\
        \textbf{Verborgenheit} & +1\\
        \textbf{Schadenspotential} & +1\\
        \textbf{Genauigkeit} & +1\\
    \end{tabular}

    \wbif{phasesix}{
        \vspace{3mm}
\faIcon{plane}\hspace{0.5em}\faIcon{microchip}\hspace{0.5em}    }{}



    \subsection*{Rotpunktvisier}

    \begin{pexample}
    Das Rotpunktvisier ermöglicht leichteres Zielen auf kurze Distanzen (<200m).
    \end{pexample}


    \begin{tabular}{@{}ll@{}}
        \textbf{\trans{price_label}:} & 1200 \\
        \textbf{Verborgenheit} & -1\\
        \textbf{Schadenspotential} & +1\\
    \end{tabular}

    \wbif{phasesix}{
        \vspace{3mm}
\faIcon{plane}\hspace{0.5em}\faIcon{microchip}\hspace{0.5em}    }{}



    \subsection*{2x Visier}


    Das 2x Visier ermöglicht ein genaues Zielen auf mittlere Distanzen, und senkt die Schwelle für kritische Treffer um 1.

    \begin{tabular}{@{}ll@{}}
        \textbf{\trans{price_label}:} & 1000 \\
        \textbf{Kritische Treffer} & -1\\
    \end{tabular}

    \wbif{phasesix}{
        \vspace{3mm}
\faIcon{flag}\hspace{0.5em}\faIcon{plane}\hspace{0.5em}\faIcon{microchip}\hspace{0.5em}\faIcon{space-shuttle}\hspace{0.5em}    }{}



    \subsection*{Armbrust-Eisen-Visier}

    \begin{pexample}
    Das Anbringen eines eisernen Visiers an der Armbrust erhöht die Treffsicherheit und den möglichen Schaden bei einem Treffer.
    \end{pexample}


    \begin{tabular}{@{}ll@{}}
        \textbf{\trans{price_label}:} & 600 \\
        \textbf{Genauigkeit} & +1\\
        \textbf{Schadenspotential} & +1\\
    \end{tabular}

    \wbif{phasesix}{
        \vspace{3mm}
\faIcon{cube}\hspace{0.5em}    }{}


    \section{Sondervorrichtung}



    \subsection*{Sturmlicht}

    \begin{pexample}
    Das Sturmlicht ist eine Lichtquelle, die den Bereich vor dem Waffenträger ausleuchtet. Das Licht bewegt sich mit der Waffe. Die Genauigkeit wird erhöht, aber der Waffenträger ist leicht zu erkennen.
    \end{pexample}


    \begin{tabular}{@{}ll@{}}
        \textbf{\trans{price_label}:} & 200 \\
        \textbf{Verborgenheit} & +2\\
        \textbf{Genauigkeit} & +1\\
    \end{tabular}

    \wbif{phasesix}{
        \vspace{3mm}
\faIcon{plane}\hspace{0.5em}\faIcon{microchip}\hspace{0.5em}    }{}



    \subsection*{Dreibein}

    \begin{pexample}
    Das Dreibein sorgt bei der Verwendung mit einem Sturmgewehr oder Maschinengewehr für eine erhebliche Reduzierung des Rückstoßes. Dafür ist das Nachladen der Waffe umständlicher.
    \end{pexample}


    \begin{tabular}{@{}ll@{}}
        \textbf{\trans{price_label}:} & 500 \\
        \textbf{Rückstosskontrolle} & +2\\
        \textbf{Verborgenheit} & +1\\
        \textbf{Nachladeaktionen} & +1\\
    \end{tabular}

    \wbif{phasesix}{
        \vspace{3mm}
\faIcon{umbrella}\hspace{0.5em}\faIcon{flag}\hspace{0.5em}\faIcon{plane}\hspace{0.5em}\faIcon{microchip}\hspace{0.5em}    }{}



    \subsection*{Schnellziehköcher}

    \begin{pexample}
    Dieser Köcher ist so konstruiert, dass ein Pfeil viel schneller auf die Sehne eines Bogens gesetzt werden kann.
    \end{pexample}


    \begin{tabular}{@{}ll@{}}
        \textbf{\trans{price_label}:} & 200 \\
        \textbf{Nachladeaktionen} & -1\\
    \end{tabular}

    \wbif{phasesix}{
        \vspace{3mm}
\faIcon{cube}\hspace{0.5em}    }{}



    \subsection*{Schnellziehschlinge}

    \begin{pexample}
    Diese Vorrichtung an der Waffe ermöglicht ein schnelles Ziehen und Schießen.
    \end{pexample}


    \begin{tabular}{@{}ll@{}}
        \textbf{\trans{price_label}:} & 200 \\
        \textbf{Vorbereitung} & -1\\
    \end{tabular}

    \wbif{phasesix}{
        \vspace{3mm}
\faIcon{cube}\hspace{0.5em}    }{}



    \subsection*{Schnellader für Revolver}

    \begin{pexample}
    Eine Vorrichtung zur Aufnahme von sechs Patronen. Damit kann ein Revolver sehr schnell geladen werden. Das Bestücken des Schnellladers dauert jedoch genauso lange wie das manuelle Laden eines Revolvers.
    \end{pexample}


    \begin{tabular}{@{}ll@{}}
        \textbf{\trans{price_label}:} & 25 \\
        \textbf{Nachladeaktionen} & -1\\
    \end{tabular}

    \wbif{phasesix}{
        \vspace{3mm}
\faIcon{umbrella}\hspace{0.5em}\faIcon{flag}\hspace{0.5em}\faIcon{plane}\hspace{0.5em}\faIcon{microchip}\hspace{0.5em}\faIcon{space-shuttle}\hspace{0.5em}    }{}



    \subsection*{Gesegnet}

    \begin{pexample}
    Die Waffe wurde von einem Priester geweiht. Auf ihr liegt der Segen einer höheren Wesenheit, sie hat besondere Fähigkeiten und wirkt stärker gegen die Mächte des Bösen.
    \end{pexample}

    Ergebnisse des Trefferwurfes, die eine 1 ergeben, dürfen einmal wiederholt werden. Die Anzahl der Treffer gegen Dämonen und Geister wird verdoppelt.

    \begin{tabular}{@{}ll@{}}
        \textbf{\trans{price_label}:} & 500 \\
        \textbf{Schadenspotential} & +1\\
    \end{tabular}

    \wbif{phasesix}{
        \vspace{3mm}
\faIcon{ankh}\hspace{0.5em}\faIcon{spider}\hspace{0.5em}    }{}


    \section{Griffe}



    \subsection*{Lederumwickelter Griff}

    \begin{pexample}
    Ein mit Leder überzogener Griff verbessert die Handhabung der Waffe und erhöht das Schadenspotential.
    \end{pexample}


    \begin{tabular}{@{}ll@{}}
        \textbf{\trans{price_label}:} & 80 \\
        \textbf{Schadenspotential} & +1\\
    \end{tabular}

    \wbif{phasesix}{
        \vspace{3mm}
\faIcon{eye}\hspace{0.5em}\faIcon{chess-rook}\hspace{0.5em}\faIcon{umbrella}\hspace{0.5em}\faIcon{flag}\hspace{0.5em}\faIcon{plane}\hspace{0.5em}\faIcon{microchip}\hspace{0.5em}    }{}



    \subsection*{Hartholzgriff}

    \begin{pexample}
    Ein Hartholzgriff verbessert die Handhabung der Waffe und erhöht das Schadenspotential und die Präzision.
    \end{pexample}


    \begin{tabular}{@{}ll@{}}
        \textbf{\trans{price_label}:} & 200 \\
        \textbf{Genauigkeit} & +1\\
        \textbf{Schadenspotential} & +1\\
    \end{tabular}

    \wbif{phasesix}{
        \vspace{3mm}
\faIcon{cube}\hspace{0.5em}    }{}


    \section{Munition}



    \subsection*{Leuchtspurmunition}

    \begin{pexample}
    Leuchtspurmunition erleichtert das Zielen auf einen Gegner und erhöht damit das Schadenspotenzial und die Genauigkeit. Dafür ist der Schütze leichter zu erkennen.
    \end{pexample}


    \begin{tabular}{@{}ll@{}}
        \textbf{\trans{price_label}:} & 50 \\
        \textbf{Verborgenheit} & +2\\
        \textbf{Schadenspotential} & +1\\
        \textbf{Genauigkeit} & +1\\
    \end{tabular}

    \wbif{phasesix}{
        \vspace{3mm}
\faIcon{flag}\hspace{0.5em}\faIcon{plane}\hspace{0.5em}\faIcon{microchip}\hspace{0.5em}\faIcon{space-shuttle}\hspace{0.5em}    }{}



    \subsection*{Gummigeschossmunition}

    \begin{pexample}
    Gummigeschosse verringern den Schaden und die Durchschlagskraft der Waffe, schocken aber den Gegner.
    \end{pexample}


    \begin{tabular}{@{}ll@{}}
        \textbf{\trans{price_label}:} & 200 \\
        \textbf{Durchschlag} & -1\\
        \textbf{Geschockt} & +2\\
        \textbf{Schadenspotential} & -4\\
    \end{tabular}

    \wbif{phasesix}{
        \vspace{3mm}
\faIcon{flag}\hspace{0.5em}\faIcon{plane}\hspace{0.5em}\faIcon{microchip}\hspace{0.5em}    }{}



    \subsection*{Flintenlaufgeschosse}

    \begin{pexample}
    Diese Munition ermöglicht es, aus einem Gewehr ein einzelnes Geschoss abzufeuern, das mehr Schaden anrichtet und die Reichweite des Gewehrs erhöht.
    \end{pexample}


    \begin{tabular}{@{}ll@{}}
        \textbf{\trans{price_label}:} & 100 \\
        \textbf{Schadenspotential} & +1\\
        \textbf{Reichweite} & +5\\
    \end{tabular}

    \wbif{phasesix}{
        \vspace{3mm}
\faIcon{cube}\hspace{0.5em}    }{}



    \subsection*{Erweitertes Magazin (Pistolen)}

    \begin{pexample}
    Das erweiterte Magazin fasst 7 Schuss mehr und kann für Pistolen verwendet werden.
    \end{pexample}


    \begin{tabular}{@{}ll@{}}
        \textbf{\trans{price_label}:} & 80 \\
        \textbf{Kapaziät} & +7\\
    \end{tabular}

    \wbif{phasesix}{
        \vspace{3mm}
\faIcon{cube}\hspace{0.5em}    }{}



    \subsection*{Erweitrertes Magazin (Maschinengewehre)}

    \begin{pexample}
    Das erweiterte Magazin fasst 20 Schuss mehr und kann für Maschinengewehre verwendet werden.
    \end{pexample}


    \begin{tabular}{@{}ll@{}}
        \textbf{\trans{price_label}:} & 150 \\
        \textbf{Kapaziät} & +20\\
    \end{tabular}

    \wbif{phasesix}{
        \vspace{3mm}
\faIcon{cube}\hspace{0.5em}    }{}



    \subsection*{Kieselsteine}

    \begin{pexample}
    Einfache Kieselsteine zur Benutzung in einer Schleuder oder Zwille.
    \end{pexample}


    \begin{tabular}{@{}ll@{}}
        \textbf{\trans{price_label}:} & 2 \\
    \end{tabular}

    \wbif{phasesix}{
        \vspace{3mm}
\faIcon{cube}\hspace{0.5em}    }{}



    \subsection*{Eisenkugeln}

    \begin{pexample}
    Eisenkugeln fügen dem Gegner mehr Schaden zu, wenn sie anstelle von Steinen in einer Schleuder verwendet werden.
    \end{pexample}


    \begin{tabular}{@{}ll@{}}
        \textbf{\trans{price_label}:} & 10 \\
        \textbf{Schadenspotential} & +2\\
    \end{tabular}

    \wbif{phasesix}{
        \vspace{3mm}
\faIcon{cube}\hspace{0.5em}    }{}



    \subsection*{Silbermunition}

    \begin{pexample}
    Aus Silber gefertigte oder mit Silber beschichtete Munition.
    \end{pexample}

    Bei einem Angriff gegen Werwölfe und Vampire dürfen alle Würfel des Angriffswurfes, die eine 1 zeigen, erneut geworfen werden.

    \begin{tabular}{@{}ll@{}}
        \textbf{\trans{price_label}:} & 50 \\
    \end{tabular}

    \wbif{phasesix}{
        \vspace{3mm}
\faIcon{spider}\hspace{0.5em}    }{}



    \subsection*{Giftpfeile}

    \begin{pexample}
    Giftpfeile haben eine spezielle Spitze, an der das Gift besonders gut haftet. Diese Pfeile verursachen Vergiftungen, die der Stärke des verwendeten Giftes entsprechen.
    \end{pexample}


    \begin{tabular}{@{}ll@{}}
        \textbf{\trans{price_label}:} & 20 \\
        \textbf{Giftkanal} & +1\\
    \end{tabular}

    \wbif{phasesix}{
        \vspace{3mm}
\faIcon{cube}\hspace{0.5em}    }{}



    \subsection*{Explosive Pfeile}

    \begin{pexample}
    Durch eine spezielle Vorrichtung an der Pfeilspitze explodieren die Pfeile beim Aufprall.
    \end{pexample}


    \begin{tabular}{@{}ll@{}}
        \textbf{\trans{price_label}:} & 700 \\
        \textbf{Flächenschaden} & +2\\
    \end{tabular}

    \wbif{phasesix}{
        \vspace{3mm}
\faIcon{cube}\hspace{0.5em}    }{}


    \section{Klinge}



    \subsection*{Aufgerauhte Klinge}

    \begin{pexample}
    Wird die Klinge einer Waffe aufgerauht, so verringert sich zwar die Durchschlagskraft der Waffe, ein Schlag verursacht jedoch stark blutende Wunden.
    \end{pexample}


    \begin{tabular}{@{}ll@{}}
        \textbf{\trans{price_label}:} & 100 \\
        \textbf{Durchschlag} & -1\\
        \textbf{Blutend} & +2\\
    \end{tabular}

    \wbif{phasesix}{
        \vspace{3mm}
\faIcon{cube}\hspace{0.5em}    }{}



    \subsection*{Gehärtete Klinge}

    \begin{pexample}
    Die gehärtete Klinge erhöht die Durchschlagskraft und das Schadenspotential der Waffe.
    \end{pexample}


    \begin{tabular}{@{}ll@{}}
        \textbf{\trans{price_label}:} & 200 \\
        \textbf{Schadenspotential} & +1\\
        \textbf{Durchschlag} & +1\\
    \end{tabular}

    \wbif{phasesix}{
        \vspace{3mm}
\faIcon{cube}\hspace{0.5em}    }{}



    \subsection*{Gravierte Klinge}

    \begin{pexample}
    Die Klinge der Waffe enthält eine besondere Gravur.
    \end{pexample}


    \begin{tabular}{@{}ll@{}}
        \textbf{\trans{price_label}:} & 100 \\
        \textbf{Schadenspotential} & +1\\
    \end{tabular}

    \wbif{phasesix}{
        \vspace{3mm}
\faIcon{cube}\hspace{0.5em}    }{}



    \subsection*{Giftkanal}

    \begin{pexample}
    Eine Kerbe zum Auftragen von Gift. Klingenwaffen können damit modifiziert werden. Vergiftet in der Stärke des verwendeten Giftes.
    \end{pexample}


    \begin{tabular}{@{}ll@{}}
        \textbf{\trans{price_label}:} & 250 \\
        \textbf{Giftkanal} & +1\\
    \end{tabular}

    \wbif{phasesix}{
        \vspace{3mm}
\faIcon{cube}\hspace{0.5em}    }{}



    \subsection*{Versilberte Klinge}

    \begin{pexample}
    Die Klinge ist versilbert und sorgt für effektivere Angriffe gegen Werwölfe und Vampire.
    \end{pexample}

    Wird die Waffe gegen Vampire oder Werwölfe eingesetzt, so werden die Treffer nach dem Trefferwurf verdoppelt.

    \begin{tabular}{@{}ll@{}}
        \textbf{\trans{price_label}:} & 550 \\
    \end{tabular}

    \wbif{phasesix}{
        \vspace{3mm}
\faIcon{spider}\hspace{0.5em}    }{}



    \subsection*{Gekrümmte Klinge}

    \begin{pexample}
    Wenn die Waffe eine gekrümmte Klinge hat, wird die Reichweite erhöht und die Wunde blutet, da die Waffe häufiger ungeschützte Körperteile erreicht.
Eine bestehende Waffe kann nicht durch einen Schmied auf eine gekrümmte Klinge umgerüstet werden, dies muss bei neuen Waffen direkt beauftragt werden.
    \end{pexample}


    \begin{tabular}{@{}ll@{}}
        \textbf{\trans{price_label}:} & 300 \\
        \textbf{Reichweite} & +1\\
        \textbf{Blutend} & +1\\
    \end{tabular}

    \wbif{phasesix}{
        \vspace{3mm}
\faIcon{cube}\hspace{0.5em}    }{}



    \subsection*{Gezackte Schneide}

    \begin{pexample}
    Eine gezackte Schneide verursacht stark blutende Wunden.
    \end{pexample}


    \begin{tabular}{@{}ll@{}}
        \textbf{\trans{price_label}:} & 400 \\
        \textbf{Blutend} & +1\\
    \end{tabular}

    \wbif{phasesix}{
        \vspace{3mm}
\faIcon{cube}\hspace{0.5em}    }{}


\end{multicols}